% Section 5.2, from lectures/L05/appendix/b_the_driver.md.

\section{The Driver, Specified}
\label{c5:app:b}

This is what to build. Everything here is prose rather than code, because the code is yours; what
the course provides is the contract, the header that declares it, and the tests that check it.

Two modules, in this order:
\begin{enumerate}
  \item \textbf{\file{qa_fifo.c}}, a byte ring buffer with no kernel API in it beyond allocation.
    Declared by \file{lectures/L05/lab/qa_fifo.h} and checked by
    \file{lectures/L05/lab/qa_fifo_kunit.c}, both of which ship with the course.
  \item \textbf{\file{qa_chardev.c}}, a character device that exposes one of those buffers at
    \file{/dev/qa_fifo}. Checked by \file{lectures/L05/lab/user/qa_fifo_test.c}, which also ships.
\end{enumerate}
Write them in that order. The first is pure logic and can be got completely right before anything
touches the kernel's device model; the second is where the kernel API is, and debugging both at once
is what makes a first driver miserable.

\subsection{The ring buffer}
\label{c5:app:b1}

The declarations are in \file{lectures/L05/lab/qa_fifo.h}, and the prose in that header is part of
the specification rather than a comment on it. Here it is, typeset straight from the repository:

\cfile{lectures/L05/lab/qa_fifo.h}

\noindent What follows is what the header cannot say.

\textbf{The struct is opaque.} \code{struct qa_fifo} is declared in the header and defined only in
your \code{.c}. Nothing outside may look inside it, which means you choose the layout, and the tests
cannot accidentally depend on your choices. It also means every access goes through a function,
which is the habit that makes Chapter~\ref{c7:ch}'s locking tractable: there is exactly one place
per operation to put a lock.

\textbf{Every function must be exported.} The KUnit suite is a separate module, so each function in
the header needs an \code{EXPORT_SYMBOL_GPL}, exactly as in Chapter~\ref{c4:ch}. Forgetting this
produces \code{Unknown symbol} at \code{insmod} time for the \emph{test} module, and the lab's
\file{run.sh} distinguishes that case from a missing KUnit framework and says which it is.

\textbf{The capacity is what you asked for.} A FIFO created with capacity 4 must hold 4 bytes, not
3. The usual ring-buffer implementation keeps head and tail indices and cannot distinguish full from
empty, so it wastes one slot; that implementation fails \code{qa_fifo_uses_its_whole_capacity}.
Keeping a \emph{level} alongside the head, rather than a tail, is one way out and there are others.

\textbf{Short is not an error.} \code{qa_fifo_put} given more than fits copies what fits and returns
that count; \code{qa_fifo_get} on a partly filled buffer copies what is there. Neither is a failure,
and the caller is expected to read the return value. This is the same contract \code{read()} and
\code{write()} have, deliberately, so that the character device can pass its arguments almost
straight through.

\textbf{Wrapping is the part that gets written wrong.} Fill, drain, and fill again: the second fill
straddles the end of the storage. An implementation that tracks positions without wrapping passes
every test in the suite up to the seventh and fails the eighth.

\subsection{The character device}
\label{c5:app:b2}

\file{qa_chardev.c} registers one device, backed by one \code{qa_fifo}.

\subsubsection{What it must provide}
\label{c5:app:b:what-it-must-provide}

\begin{narrowtable}{@{}lll@{}}
  \thead{Member} & \thead{Required} & \thead{Behaviour} \\
  \midrule
  \code{.owner} & \textbf{Yes} & \code{THIS_MODULE}. See \S\ref{c5:app:a3} \\
  \code{.open} & Yes & Succeeds. Marks the file as a stream \\
  \code{.read} & Yes & Below \\
  \code{.write} & Yes & Below \\
  \code{.release} & No & There is no per-open state to release \\
\end{narrowtable}

\subsubsection{Module parameter}
\label{c5:app:b:module-parameter}

One parameter, \code{capacity}, an \code{int}, default \textbf{256}, readable but not writable
through sysfs. It is the capacity of the FIFO created at load time. The userspace tester does not
assume a value; it measures what it can write and reports it, so changing the default does not break
the test.

\subsubsection{\code{open}}
\label{c5:app:b:open}

Succeed, and mark the file as a stream. \code{stream_open()} does exactly this and is the one line
the function needs. Note that \code{no_llseek}, which most older material uses for this, \textbf{was
removed and does not exist in 6.12}.

A FIFO has no position, and a driver that silently accepts an offset is lying about what it is. Two
different things refuse one, and it is worth knowing which does what. Since 6.0 the kernel refuses
\code{lseek} with \code{-ESPIPE} on any file whose \code{file_operations} has no \code{.llseek},
which yours does not; the userspace tester checks that, and it passes with or without
\code{stream_open}. \textbf{\code{stream_open} is what refuses the rest.} It makes \code{pread} and
\code{pwrite} fail with \code{-ESPIPE} too, where without it they would succeed and read the FIFO as
though the offset meant something, and it tells the VFS there is no file position at all, so the
\code{pos} your \code{read} and \code{write} are handed is \code{NULL} and must not be touched.

\subsubsection{\code{read}}
\label{c5:app:b:read}

Given a userspace buffer and a count:
\begin{itemize}
  \item A \code{count} of 0 returns 0. There is nothing to do and it is not an error.
  \item \textbf{If the FIFO is empty, return \code{-EAGAIN}.} Not 0.
  \item Otherwise take up to \code{count} bytes out of the FIFO and copy them to userspace, and
    return how many. Fewer than \code{count} is normal.
  \item If the copy fails, return \code{-EFAULT}.
\end{itemize}
You cannot copy straight from the FIFO to the user pointer, because \S\ref{c5:app:a5} says you may
not touch that pointer directly. So there is a bounce buffer in the middle: allocate, take from the
FIFO into it, \code{copy_to_user} out of it, free. Cap the allocation at something sensible; a
program is allowed to ask for a gigabyte and you are not obliged to allocate one.

\textbf{The \code{-EAGAIN} is the whole chapter.} Returning 0 means end of file, and every reader
believes it. Chapter~\ref{c9:ch} replaces this with blocking, at which point \code{-EAGAIN} becomes
the answer only for a caller that opened with \code{O_NONBLOCK}.

\subsubsection{\code{write}}
\label{c5:app:b:write}

The mirror image:
\begin{itemize}
  \item A \code{count} of 0 returns 0.
  \item \textbf{If the FIFO is full, return \code{-ENOSPC}.}
  \item Otherwise copy up to \code{count} bytes in from userspace and return how many were accepted.
  \item If the copy fails, return \code{-EFAULT}.
\end{itemize}
Note the asymmetry with \code{read}: a full FIFO is \code{-ENOSPC} and an empty one is
\code{-EAGAIN}. Both say ``not now'', but a reader can reasonably wait for data to arrive, whereas a
writer facing a buffer nobody is draining has a different problem. The codes are different because
userspace does different things with them.

\subsubsection{Init, in the right order}
\label{c5:app:b:init-in-the-right-order}

Five things to acquire, and the order is a constraint rather than a preference:
\begin{enumerate}
  \item Create the FIFO.
  \item Allocate a device number \textbf{dynamically}, with \code{alloc_chrdev_region}.
  \item \code{cdev_init} and \code{cdev_add}.
  \item \code{class_create}.
  \item \code{device_create}, which is what makes \file{/dev/qa_fifo} appear.
\end{enumerate}
\textbf{Everything your callbacks touch must exist before step 3.} \code{cdev_add} makes the device
live; from that instant a \code{read()} can arrive, and it will arrive before your init function has
returned if anything is watching \file{/dev}. Creating the FIFO after adding the cdev is a race you
will not reproduce on demand and will hit eventually.

Print the major and minor you were given at the end, with \code{pr_info}. You will want it, and it
is different on every boot.

\subsubsection{Exit, in reverse}
\label{c5:app:b:exit-in-reverse}

Undo all five, in the opposite order, and free the FIFO last. Any error path in init must unwind
only what it actually acquired; see \S\ref{c5:app:a7} for the \code{goto} ladder that expresses
this.

\subsection{What the tests check, and what they do not}
\label{c5:app:b3}

\textbf{The KUnit suite} (\file{lectures/L05/lab/qa_fifo_kunit.c}) is ten cases against the ring
buffer, run inside the kernel. It says so at the top, and it is worth repeating here: \textbf{it
checks nothing about concurrency.} Every case runs on one thread. A ring buffer that passes all ten
and is destroyed by two writers passes all ten. That is Chapter~\ref{c7:ch}.

\textbf{The userspace tester} (\file{lectures/L05/lab/user/qa_fifo_test.c}) is a statically linked
ARM64 program that opens the node and makes one system call at a time. It exists as a program rather
than as shell because several of the things a driver must get right are invisible from a shell: a
\code{read} returning 0 and a \code{read} returning \code{-EAGAIN} both produce no output from
\code{cat}, and only a program can look at \code{errno}.

Between them they do \textbf{not} check: concurrency, what happens under memory pressure, whether
your error paths unwind correctly, or whether you leak the FIFO on an init failure. Reading the
tests to find out what they are blind to is a habit worth forming, and one of the exercises asks you
to do exactly that.

\subsection{Working order}
\label{c5:app:b4}

\begin{enumerate}
  \item Write \file{qa_fifo.c} and get the KUnit suite to \textbf{10 of 10}. Nothing below matters
    until it does.
  \item Write \file{qa_chardev.c} with \code{open}, \code{read} and \code{write}, and load it. Check
    \file{/proc/devices} for your major and \file{/dev/qa_fifo} for the node.
  \item Run \file{./qa_fifo_test} on the target by hand before running \code{make test}; its output
    is more informative than the harness's summary.
  \item \code{make test L=L05} reports \textbf{PASSED}, with eighteen checks.
\end{enumerate}
If the KUnit module refuses to load with \code{Unknown symbol qa_fifo_create}, you have not exported
the functions. If it refuses with \code{Unknown symbol kunit_...}, the framework is not loaded and
\file{run.sh} will say so.
